\subsubsection{A perturbative weighted quadrature-domain theorem}\label{app:cherednichenko-type}

We now prove a local realization statement for weighted quadrature domains
near a reference \(C^2\) domain. The result is a tailored, perturbative
version of the classical setup of Cherednichenko \cite{Cherednichenko},
adapted to the form needed in the proof of Theorem~\ref{thm:perturb}.
Conceptually, the statement may also be viewed as a quantitative
implicit-function counterpart of the variational characterization of
quadrature domains studied by Gustafsson \cite{Gustafsson},
Gustafsson--Shahgholian \cite{Gustafsson-Shahgholian},
Shahgholian \cite{Shahgholian}, Karp \cite{Karp}, and
Aharonov--Schiffer--Zalcman \cite{Aharonov-Schiffer-Zalcman}.

\medskip
\noindent\emph{Setup.}
Fix a positive weight \(\mu\in C^2(\overline{U_2})\), where
\(U_1\subset U_0\subset U_2\) are bounded \(C^2\) domains with
\(\overline{U_1}\subset U_0\) and \(\overline{U_0}\subset U_2\).
For a bounded measurable domain \(U\subset U_2\), define its weighted
Cauchy transform by
\[
g_U(z):=\int_U \frac{\mu(w)}{z-w}\,dA(w),
\qquad z\in \C\setminus \overline U.
\]
The function \(g_U\) is holomorphic on \(\C\setminus \overline U\), continuous up to
\(\partial U\), and, in the sense of distributions on \(\C\),
\[
\partial_{\overline z}\, g_U \;=\; \pi\,\mu\,\mathbf{1}_U.
\]
At infinity it admits the expansion
\[
g_U(z)=\frac{1}{z}\int_U \mu(w)\,dA(w)+O(|z|^{-2}),
\]
so in particular \(\lim_{z\to\infty}g_U(z)=0\) and
\(\lim_{z\to\infty}zg_U(z)=\int_U\mu\,dA>0\).

As before, every domain \(U\) sufficiently \(C^2\)-close to \(U_0\) may be
represented as a normal graph over \(\partial U_0\),
\[
\partial U
=
\{x+\varphi(x)\,\nu(x):x\in \partial U_0\},
\]
with \(\varphi\in C^2(\partial U_0)\) of small \(C^2\) norm, where \(\nu\)
denotes the outward unit normal to \(\partial U_0\). We denote the resulting
domain by \(U_\varphi\) and write \(\Phi_\varphi(x):=x+\varphi(x)\nu(x)\) for
the boundary diffeomorphism, extended smoothly to a tubular neighborhood of
\(\partial U_0\) (and by the identity outside) so as to define a
diffeomorphism of \(\overline{U_2}\) onto itself.

\begin{proposition}[\cite{Cherednichenko}]\label{prop:cherednichenko-tailored}
Let \(U_1\subset U_0\subset U_2\) be bounded \(C^2\) domains, and let
\(\mu\in C^2(\overline{U_2})\) satisfy
\[
\mu>0 \qquad \text{on }\overline{U_2}.
\]
Then there exist \(\rho>0\) and \(\varepsilon>0\) with the following property.

Suppose \(g\) is holomorphic on \(\C\setminus \overline{U_1}\), continuous up
to \(\partial U_1\), and satisfies
\begin{equation}\label{eq:appendix-g-hypotheses}
\lim_{z\to\infty} g(z)=0,
\qquad
\lim_{z\to\infty} z\,g(z)=c_0>0,
\end{equation}
together with
\[
\sup_{z\in \partial U_1}|g(z)-g_{U_0}(z)|<\varepsilon.
\]
Then there exists a unique function
\(\varphi\in C^2(\partial U_0)\), with \(\|\varphi\|_{C^2(\partial U_0)}<\rho\),
such that the associated domain \(U_\varphi\) satisfies
\[
U_1\subset U_\varphi\subset U_2
\]
and
\[
g_{U_\varphi}(z)=g(z)
\qquad \text{for all } z\in \C\setminus \overline{U_\varphi}.
\]
Moreover, the assignment
\[
g\mapsto \varphi(g)
\]
is \(C^1\) with respect to the \(C^2(\partial U_1)\) topology on the trace
data and the \(C^2(\partial U_0)\) topology on the boundary graph.
\end{proposition}

\begin{proof}
The proof reduces matters to a boundary equation on \(\partial U_1\) and
applies the implicit function theorem in suitable Banach spaces.

\medskip
\noindent\emph{Step 1: reduction to a boundary matching problem.}
Both \(g\) and \(g_{U_\varphi}\) are holomorphic on
\(\C\setminus \overline{U_\varphi}\) and vanish at infinity. As long as
\(U_1\subset U_\varphi\), the curve \(\partial U_1\) lies in the (connected)
unbounded component of \(\C\setminus \overline{U_\varphi}\). Suppose
\begin{equation}\label{eq:boundary-equation-U1}
g_{U_\varphi}(z)\;=\;g(z)
\qquad \text{for all }z\in \partial U_1.
\end{equation}
The difference \(g_{U_\varphi}-g\) is then holomorphic on
\(\C\setminus \overline{U_\varphi}\), continuous up to the boundary, vanishes
at infinity, and vanishes on the closed curve \(\partial U_1\). Since the
zero set of a non-trivial holomorphic function on a connected open set is
discrete, vanishing on the one-dimensional set \(\partial U_1\) forces
\(g_{U_\varphi}\equiv g\) on the unbounded component of
\(\C\setminus \overline{U_\varphi}\). As both functions are holomorphic on
the full open set \(\C\setminus \overline{U_\varphi}\), the identity theorem
extends the equality to that whole open set.

Thus it suffices to solve \eqref{eq:boundary-equation-U1}. We define
\[
\mathcal{G}:\;\varphi\;\longmapsto\;g_{U_\varphi}\big|_{\partial U_1},
\]
and view \(\mathcal{G}\) as a map between Banach spaces
\[
\mathcal{G}\colon\;
\mathcal{O}\subset C^2(\partial U_0)\;\longrightarrow\;C^2(\partial U_1),
\]
defined on a small \(C^2\) neighborhood \(\mathcal{O}\) of \(0\), small
enough that \(\overline{U_1}\subset U_\varphi\subset U_2\) for all
\(\varphi\in \mathcal{O}\). Note that for such \(\varphi\), \(\partial U_1\)
stays a positive distance from \(\overline{U_\varphi}\), so the integral
defining \(g_{U_\varphi}\) is non-singular for \(z\in \partial U_1\).

\medskip
\noindent\emph{Step 2: \(C^1\) dependence of \(\mathcal{G}\) on \(\varphi\).}
Performing the change of variables \(w=\Phi_\varphi(w')\) in the bulk integral
yields
\[
g_{U_\varphi}(z)
=
\int_{U_0}
\frac{\mu(\Phi_\varphi(w'))\,J_\varphi(w')}{z-\Phi_\varphi(w')}\,dA(w'),
\qquad z\in \partial U_1,
\]
where \(J_\varphi=|\det D\Phi_\varphi|\). Because \(z\in \partial U_1\) keeps
a uniform positive distance \(d_0>0\) from
\(\overline{U_0}\) (and hence from \(\overline{U_\varphi}\) for all
\(\varphi\) in a sufficiently small neighborhood of zero), the integrand is a
smooth function of \((\varphi,w')\) and the differentiation in \(\varphi\)
under the integral sign is fully justified by dominated convergence.
Combined with the smoothness of \(\mu\) and of the dependence
\(\varphi\mapsto \Phi_\varphi\), this gives that \(\mathcal{G}\) is of class
\(C^1\) (in fact \(C^k\) whenever \(\mu\in C^k\)) from
\(\mathcal{O}\subset C^2(\partial U_0)\) to \(C^2(\partial U_1)\).

\medskip
\noindent\emph{Step 3: Hadamard shape derivative.}
We compute \(D\mathcal{G}(0)\). Let \(\psi\in C^2(\partial U_0)\) and set
\(\varphi=t\psi\) for small \(t\in \R\). Writing
\(U_{t\psi}\setminus U_0\) as the set swept out by the normal flow with
speed \(\psi\) in time \(t\), one has, for \(z\in \partial U_1\) (so that
\(z\notin \overline{U_{t\psi}}\) for \(|t|\) small),
\[
g_{U_{t\psi}}(z)-g_{U_0}(z)
=
\int_{U_{t\psi}\triangle U_0}\frac{\pm \mu(w)}{z-w}\,dA(w),
\]
with the sign \(+\) (resp.\ \(-\)) on \(U_{t\psi}\setminus U_0\) where
\(\psi>0\) (resp.\ on \(U_0\setminus U_{t\psi}\) where \(\psi<0\)).
Parametrizing the symmetric difference as a tubular shell of width
\(t\psi(x)\) over \(\partial U_0\) and using the surface area element
\(d\sigma\) on \(\partial U_0\), Fubini gives
\[
g_{U_{t\psi}}(z)-g_{U_0}(z)
=
\int_{\partial U_0}
\Bigl(\int_0^{t\psi(x)}\frac{\mu(x+s\nu(x))}{z-x-s\nu(x)}\,ds\Bigr)\,d\sigma(x),
\]
which after dividing by \(t\) and letting \(t\to 0\) yields the standard
Hadamard variation formula
\begin{equation}\label{eq:appendix-shape-derivative-cauchy}
D\mathcal{G}(0)[\psi](z)
=
\int_{\partial U_0}
\frac{\mu(w)\,\psi(w)}{z-w}\,d\sigma(w),
\qquad z\in \partial U_1.
\end{equation}
Thus the linearization is the (single-layer) Cauchy transform on
\(\partial U_0\) of the boundary density \(\mu\psi\), restricted to
\(\partial U_1\). We denote it by \(L\),
\[
L\psi(z):=\int_{\partial U_0}\frac{\mu(w)\,\psi(w)}{z-w}\,d\sigma(w),
\qquad z\in \partial U_1.
\]

\medskip
\noindent\emph{Step 4: invertibility of the linearization.}
We show that \(L\) is an isomorphism between \(C^2(\partial U_0)\) and the
natural trace space
\[
\mathcal{T}:=\Bigl\{H\big|_{\partial U_1}:H\text{ holomorphic on }
\C\setminus \overline{U_1},\ H(\infty)=0\Bigr\}\cap C^2(\partial U_1),
\]
equipped with the \(C^2(\partial U_1)\) norm.

\smallskip
\emph{Injectivity.} Suppose \(L\psi\equiv 0\) on \(\partial U_1\). The
Cauchy-type transform
\[
F(z):=\int_{\partial U_0}\frac{\mu(w)\,\psi(w)}{z-w}\,d\sigma(w)
\]
is holomorphic on \(\C\setminus \partial U_0\) and vanishes at infinity.
By hypothesis \(F\equiv 0\) on \(\partial U_1\); the same identity-theorem
argument as in Step~1 (applied on the connected open set
\(\C\setminus \overline{U_0}\)) shows that
\(F\equiv 0\) on \(\C\setminus \overline{U_0}\). The Sokhotski--Plemelj
jump relations for the Cauchy integral on the \(C^2\) curve
\(\partial U_0\) \cite[Ch.~III]{Garnett} yield, for every
\(\zeta\in \partial U_0\) at which \(\mu\psi\) is continuous,
\[
F^{\mathrm{int}}(\zeta)-F^{\mathrm{ext}}(\zeta)
\;=\;c_0\,\mu(\zeta)\psi(\zeta),
\]
where \(F^{\mathrm{int}}\) and \(F^{\mathrm{ext}}\) denote the
non-tangential limits from the interior and the exterior of \(U_0\)
respectively, and \(c_0\) is a non-zero constant depending only on the
chosen orientation conventions for the kernel \(1/(z-w)\,d\sigma\)
(in the standard convention with \(\partial U_0\) oriented counterclockwise
and \(d\sigma\) the arclength measure, \(c_0\) equals a non-zero scalar
multiple of the integrand of \(dw/d\sigma\) on \(\partial U_0\); the
precise value is irrelevant here). Since \(F^{\mathrm{ext}}\equiv 0\) and
\(\mu>0\) on \(\partial U_0\), this forces \(\psi\equiv 0\). Hence \(L\)
is injective.

\smallskip
\emph{Surjectivity.} Let \(h\in \mathcal{T}\) and let \(H\) denote its
unique holomorphic extension to \(\C\setminus \overline{U_1}\) vanishing at
infinity. The annular region \(A:=U_0\setminus \overline{U_1}\) is
contained in \(\C\setminus \overline{U_1}\), so \(H\) is holomorphic in a
neighborhood of \(\overline{A}\); in particular the trace
\(H|_{\partial U_0}\) is well defined, lies in \(C^2(\partial U_0)\), and
depends continuously on \(h\) in the \(C^2\) topology by the Cauchy integral
formula on \(\partial U_1\) (interior holomorphic regularity inside the
annulus). We now seek a density \(\sigma_0\in C^2(\partial U_0)\) such that
\[
H(z)\;=\;\int_{\partial U_0}\frac{\sigma_0(w)}{z-w}\,d\sigma(w)
\qquad \text{for all }z\in \C\setminus \overline{U_0}.
\]
Such a density is unique: if two densities \(\sigma_0\) and
\(\sigma_0'\) produced the same Cauchy integral on
\(\C\setminus \overline{U_0}\), then their difference \(\delta:=\sigma_0
-\sigma_0'\) generates a Cauchy integral \(F\) with exterior limit
\(F^{\mathrm{ext}}\equiv 0\). The standard Sokhotski--Plemelj jump
relations \cite[Ch.~III]{Garnett} yield
\(F^{\mathrm{int}}-F^{\mathrm{ext}}=c_0\,\delta\) at every continuity
point of \(\delta\), with the same non-zero kernel constant \(c_0\) as
above. With \(F^{\mathrm{ext}}\equiv 0\), the function \(F\) extends
holomorphically across \(\partial U_0\) to the whole plane (and remains
holomorphic) only if \(F^{\mathrm{int}}\equiv 0\) as well; combined with
the vanishing of \(F\) at infinity this would force \(F\equiv 0\) and
hence \(\delta\equiv 0\). One then verifies the consistency by noting
that the \(C^2\) regularity of the curve and density guarantees that the
non-tangential limits exist pointwise and the jump identity holds in the
classical sense.

Existence of \(\sigma_0\in C^2(\partial U_0)\) is the classical exterior
Cauchy inversion problem on the \(C^2\) curve \(\partial U_0\). The
boundary integral operator
\[
T\sigma_0(\zeta)
\;:=\;\mathrm{p.v.}\!\int_{\partial U_0}
\frac{\sigma_0(w)}{\zeta-w}\,d\sigma(w),\qquad \zeta\in\partial U_0,
\]
is a classical singular integral operator on the \(C^2\) curve, bounded on
\(C^2(\partial U_0)\) by the standard mapping properties of the Cauchy
transform on smooth curves \cite[Ch.~III]{Garnett}. Combining \(T\) with
the Sokhotski--Plemelj jump identity gives, for the exterior boundary
value of the Cauchy integral with density \(\sigma_0\),
\[
F^{\mathrm{ext}}(\zeta)\;=\;
T\sigma_0(\zeta)-\tfrac{1}{2}\,c_0\,\sigma_0(\zeta),
\]
so that the exterior matching equation
\(F^{\mathrm{ext}}=H|_{\partial U_0}\) becomes the integral equation
\[
T\sigma_0-\tfrac{1}{2}\,c_0\,\sigma_0
\;=\;H\big|_{\partial U_0}
\qquad \text{on }\partial U_0.
\]
The associated Fredholm theory (with \(T\) a compact perturbation of a
multiple of the identity in the appropriate scale of H\"older spaces, see
\cite[Ch.~III]{Garnett}) combined with the uniqueness statement just
established yields, by the Fredholm alternative, a unique solution
\(\sigma_0\in C^2(\partial U_0)\) depending continuously on
\(H|_{\partial U_0}\), and hence on \(h\), in the \(C^2\) topology.

Setting \(\psi:=\sigma_0/\mu\in C^2(\partial U_0)\) (which is admissible
since \(\mu\geq c>0\) on \(\overline{U_2}\) and \(\mu\in C^2\)), we obtain
\[
L\psi(z)
=\int_{\partial U_0}\frac{\mu(w)\psi(w)}{z-w}\,d\sigma(w)
=\int_{\partial U_0}\frac{\sigma_0(w)}{z-w}\,d\sigma(w)
=H(z)
\qquad \text{for }z\in \partial U_1.
\]
Hence \(L\) is surjective onto \(\mathcal{T}\). Combining injectivity and
surjectivity, \(L\) is a continuous bijection between the Banach spaces
\(C^2(\partial U_0)\) and \(\mathcal{T}\); the open mapping theorem then
guarantees that \(L^{-1}\colon \mathcal{T}\to C^2(\partial U_0)\) is
bounded.

\medskip
\noindent\emph{Step 5: implicit function theorem.}
We now have all ingredients to invoke the implicit function theorem
between the Banach spaces
\[
X:=C^2(\partial U_0),\qquad
Y:=\bigl\{H\big|_{\partial U_1}:H\text{ holom. on }
\C\setminus \overline{U_1},\ H(\infty)=0\bigr\}\cap C^2(\partial U_1)
=\mathcal{T},
\]
both equipped with the appropriate \(C^2\) norms. By Step~2, the map
\(\mathcal{G}\colon\mathcal{O}\to Y\) is \(C^1\) on a neighborhood
\(\mathcal{O}\subset X\) of \(0\), and by Step~4 its differential at
\(0\) is the Banach isomorphism \(D\mathcal{G}(0)=L:X\to Y\). The
implicit function theorem (in Banach spaces) therefore furnishes
\(\rho_0>0\) and \(\varepsilon_0>0\) such that, for every
\(\widetilde h\in Y\) with
\(\|\widetilde h-\mathcal{G}(0)\|_{C^2(\partial U_1)}<\varepsilon_0\),
there is a unique \(\varphi\in X\) with
\(\|\varphi\|_{C^2(\partial U_0)}<\rho_0\) solving
\(\mathcal{G}(\varphi)=\widetilde h\), and the assignment
\(\widetilde h\mapsto \varphi\) is of class \(C^1\) from a
neighborhood of \(\mathcal{G}(0)=g_{U_0}|_{\partial U_1}\) in \(Y\) to
\(X\).

To match the hypotheses of the proposition, we observe that any \(g\)
holomorphic on \(\C\setminus \overline{U_1}\), continuous up to
\(\partial U_1\), and satisfying \eqref{eq:appendix-g-hypotheses}, with
\[
\sup_{z\in\partial U_1}|g(z)-g_{U_0}(z)|<\varepsilon,
\]
is in particular bounded on \(\partial U_1\); since \(g-g_{U_0}\) is
holomorphic on \(\C\setminus \overline{U_1}\) and vanishes at infinity,
the Cauchy integral formula on a slightly larger smooth contour
surrounding \(\partial U_1\) inside \(U_0\setminus \overline{U_1}\)
upgrades the sup-norm bound on \(\partial U_1\) to a \(C^2\)-norm bound
on \(\partial U_1\), namely
\[
\|g-g_{U_0}\|_{C^2(\partial U_1)}\;\leq\;C\,
\sup_{z\in\partial U_1}|g(z)-g_{U_0}(z)|,
\]
with a constant \(C\) depending only on \(U_0\), \(U_1\), and \(\mu\).
Choosing \(\varepsilon\) so that \(C\varepsilon<\varepsilon_0\), and
setting \(\rho:=\rho_0\), the IFT yields the unique
\(\varphi\in C^2(\partial U_0)\) with the asserted properties; in
particular \(g\mapsto\varphi(g)\) is \(C^1\) with respect to the
\(C^2(\partial U_1)\) norm on the trace data, as claimed.

By Step~1, the boundary equality \(g_{U_\varphi}=g\) on \(\partial U_1\)
extends to all of \(\C\setminus \overline{U_\varphi}\) by holomorphic
continuation.

Finally, the geometric inclusion
\(U_1\subset U_\varphi\subset U_2\) follows from a uniform tubular
neighborhood argument. Set
\[
d_1:=\mathrm{dist}(\partial U_1,\partial U_0)>0,\qquad
d_2:=\mathrm{dist}(\partial U_0,\partial U_2)>0,
\]
which are strictly positive by the standing assumption
\(\overline{U_1}\subset U_0\) and \(\overline{U_0}\subset U_2\). For
\(\varphi\in C^2(\partial U_0)\) with
\(\|\varphi\|_{C^0(\partial U_0)}<\min(d_1,d_2)\), each boundary point of
\(U_\varphi\) lies within distance \(\|\varphi\|_{C^0}\) of \(\partial U_0\),
so \(\partial U_\varphi\) is contained in the open tubular neighborhood
\(\{x:\mathrm{dist}(x,\partial U_0)<\min(d_1,d_2)\}\), which is in turn
contained in \(U_2\setminus \overline{U_1}\). Combined with the fact that
\(U_\varphi\) is a normal-graph perturbation of \(U_0\) (so its topological
position relative to \(\partial U_0\) is preserved for small \(\varphi\)),
this forces \(U_1\subset U_\varphi\subset U_2\). Since the inclusion
\(C^2(\partial U_0)\hookrightarrow C^0(\partial U_0)\) is continuous,
shrinking \(\rho\) if necessary ensures the smallness condition above.
This proves the proposition.
\end{proof}

\begin{remark}
The only delicate point in
Proposition~\ref{prop:cherednichenko-tailored} is the invertibility of the
linearized map. The mechanism is that the Cauchy transform of a boundary
density is determined by its trace on any inner contour (by holomorphic
continuation through the annulus), and the Sokhotski--Plemelj jump relation
then recovers the density itself. Positivity of \(\mu\) is what allows one
to pass from \(\mu\psi\) back to \(\psi\); the same proof would fail at any
point where \(\mu\) vanishes.
\end{remark}

\medskip

We postpone the proof of Theorem~\ref{thm:eig-loc-U} until after the
perturbative inverse theorem. The same inverse/free-boundary mechanism used
to construct domains for prescribed eigenfunctions also recovers the
classical Abreu--D\"orfler disk theorem in the simply connected setting; see
Section~\ref{sec:classical-rigidity} below.

